%-----------------------------------------------------------------------------
%
%	preface.tex Document preface part
%
%	INCLUDE FILE FOR LaTeX2e DOCUMENT
%
%	AUTHOR: Ari Potkonen /JARVENPAA/ Sat Dec 09 2023
%-----------------------------------------------------------------------------
%        1         2         3         4         5         6         7
%23456789012345678901234567890123456789012345678901234567890123456789012345678
%-BEGIN OF INCLUDE FILE-------------------------------------------------------
PREFACE
\label{preface}
\index{preface}
\vskip\baselineskip
\addcontentsline{toc}{chapter}{Preface}
Personally I have been missing clear coupling and simplicity between
electromagnetism and quantum mechanics modeling, which is fuzzy area,
overloaded with the terms. Peoples have done good work, and existing
models are usable for the practical life. Still there is need to go
back in time and lift again obvious possible connection between
electromagnetism and quantum physics, if there is any possibility to
get any clearance to our relativistic world. Now we use lazy mathematician
method: Guess the correct answer and let the others to proof that your and
other historic proposals were right. Here correct answer is found with the
simple "Duck test": "If it walks like a duck, looks like a duck, and quacks
like a duck, it's a duck". In practice I state that there is only
electromagnetism and space in our relativistic world, noting else, nothing
less. Electromagnetism itself is kind of progressing phenomenon, supporting
causality and offers us as humans possibility to define time which is handy
for modeling and tracking this causality, but time definition itself is
also based to electromagnetism and space, the major ingredients just
proposed for our generally relativistic world.
I really do not go to mathematic models, because then we overload ourselves
and loose our target under mass of details. I try to explain needed change
in our thinking in simplistic way, referring things we know from general
relativity, and keeping in mind electromagnetic space nonlinearity; high
field density -- bending electromagnetic space -- discrete curled wave
form -- both handed forms -- behave as mass -- energy dwell -- gravity.

Hopefully you'll get some ideas from here for your professional life!
\vskip\baselineskip
Sincerely yours Ari Potkonen
\vspace*{\fill}
%-END OF INCLUDE FILE---------------------------------------------------------
